\section{Normalization and the prime sum}
\label{s5-normalization}

We keep $K=40n$, $N=3n$, $h=37n$, $\alpha=3/40$, and
$\lambda=37/40$.  In particular, the scalar $S_K$ is unchanged.
The purpose of normalization is to turn the local determinant estimates
into an integral polynomial while recording the entire cost, including
positive valuations which permit division by a prime power.

\paragraph{[Conditional] The normalization exponents.}
The analogue of source (5.1)--(5.3), for fixed $M\geq40$ and
$K\geq200M^2$, is
\begin{equation}
 L_{9,p}(K,M)=
 \begin{cases}
 -10h\lfloor\log_p(5K)\rfloor-hv_p(40320),&pM\leq K,\\
 v_p(S_K)+\gamma^{\mathrm{in}}_{9,p},&pM>K,\ 3p\leq K,\\
 v_p(S_K)+\gamma^{\mathrm{out}}_{9,p},&3p>K.
 \end{cases}
 \label{s5-local-exponents}
\end{equation}
Here the two $\gamma$'s are the explicit weight sums and outer formulas
of Section~4, with $\gamma^{\mathrm{out}}_{9,p}=0$ for $p>K$.  Set
\begin{equation}
 m_{9,K,M}=\prod_{p\leq2h}p^{-L_{9,p}(K,M)},\qquad
 Q_{9,K,M}=m_{9,K,M}F_{9,K}.
 \label{s5-normalizing-product}
\end{equation}
No positive-part operation is applied to $-L_{9,p}$: the product is a
positive rational number, not necessarily an integer.  Legendre's formula gives
\begin{align}
 v_p(S_K)={}&2h\sum_{a\geq1}\lfloor K/p^a\rfloor
 -12h\sum_{a\geq1}\lfloor N/p^a\rfloor\notag\\
 &-2\sum_{i=1}^{h-1}\sum_{a\geq1}\lfloor2i/p^a\rfloor
 +(h-1)v_p(4).
 \label{s5-legendre}
\end{align}
Subject to the local estimates identified in Sections~3--4,
$Q_{9,K,M}\in\mathbb Z[X]$, the analogue of Proposition~5.1.
Indeed, these estimates handle $p\leq2h$; for $p>2h$ every factorial
in $S_K$ is a unit and the determinant is integral.  The analytic
completion step needed for the inner estimate remains a stated
dependency, rather than a theorem established by this normalization.

\subsection{The inner limiting function}
\label{s5-inner}

\paragraph{[Derived] The change from $5$ to $9$.}
For $x\geq3$ and $0<z<1/2$, put
\[
 \ell(x,z)=\lfloor x-z\rfloor+\lfloor x+z\rfloor+1,
 \qquad b(x,z)=3\ell(\alpha x,z).
\]
With $H=23/20$, define
$T=\lfloor2Hx\rfloor$, $s=Hx-T/2$, $q=\lfloor2x\rfloor$,
and $n_+=(2x-q)/2$.  The limiting weight sum is
\begin{align}
 \Gamma_9(x)={}&\int_0^{1/2}(T-b(x,z))
       (T+b(x,z)-\ell(x,z)-9)\,dz\notag\\
 &+s(2T-q-9)+(s-n_+)_+.
 \label{s5-gamma9}
\end{align}
Thus the constant $5$ in source (5.4) becomes $9$.  The allocation of
rows, including the last positive-part term, stays the same.  Since
$\int_0^{1/2}b(x,z)\,dz=3\alpha x$ and $H-3\alpha=\lambda$,
subtracting the source expression gives exactly
\[
 \Gamma_9(x)-\Gamma_5(x)
 =-4\left(T/2-3\alpha x+s\right)=-4\lambda x.
\]
The scalar still has limiting function
\begin{gather}
 J(u)=mu-\frac{m(m+1)}4,\quad m=\lfloor2u\rfloor,\qquad
 \mathcal N(x)=2\lambda x\lfloor x\rfloor
 -12\lambda x\lfloor\alpha x\rfloor-2J(\lambda x),\notag\\
 R_9(x)=-\Gamma_9(x)-\mathcal N(x)=R_5(x)+4\lambda x.
 \label{s5-r9-shift}
\end{gather}
Consequently, the analogue of (5.7) is
\[
 \gamma^{\mathrm{in}}_{9,p}=p\Gamma_9(K/p)+O_M(1),\qquad
 v_p(S_K)=p\mathcal N(K/p)+O_M(1).
\]
The finitely many reserved zero-block rows affect only the error term.
This is an algebraic comparison of the certificates; applying them to
the actual determinant still needs the local input.

\subsection{The outer limiting function}
\label{s5-outer}

\paragraph{[Derived] A vanishing rank correction.}
Write $y=p/K$.  In this range $y>1/3$, so
$(1+4\alpha-4y)_+=0$.  More concretely,
$r_{9,p}=\max(0,K+4N-4p+4)=0$ for all sufficiently large $K$.
The new raw determinant cost is therefore
\begin{equation}
 C_9(y)=
 \begin{cases}
 11(1-y)-12\alpha-6\min(\alpha,1-2y)
       -6(1+\alpha-3y)_+,&1/3<y<1/2,\\
 11(1-y)-6\min(\alpha,1-y)
       -6(1+\alpha-2y)_+,&1/2<y<1,\\
 0,&y>1.
 \end{cases}
 \label{s5-c9}
\end{equation}
For example, when $K\geq2p$, the limits of $v/K$ and $u/K$ are
$1-2y$ and $(1+\alpha-3y)_+$; inserting these in Section~4 gives
the first line.  This replaces both $R_0-d$ in source (5.8)--(5.9).
There is no additional limiting rank subtraction.

The scalar contribution is unchanged:
\[
 S(y)=-2\lambda\lfloor1/y\rfloor+
              \sum_{j=1}^{5}(2\lambda-jy)_+.
\]
Exact integration, detailed in Appendix~B, gives
\begin{gather}
 \int_{1/3}^{1}C_9(y)\,dy=\frac{14441}{7200},\qquad
 \int_{1/3}^{2\lambda}S(y)\,dy=\frac{1851}{16000},\notag\\
 I_{\mathrm{out},9}=\int_{1/3}^{2\lambda}(C_9(y)+S(y))\,dy
                   =\frac{305479}{144000}.
 \label{s5-iout9}
\end{gather}
Relative to the source's $I_{\mathrm{out},5}=127751/96000$, the
increase is $45541/57600$.  This calculation includes the improvement
in polynomial-moment integrality; it does not simply multiply every
old local loss by a common factor.

\paragraph{[Conditional] Summing over primes.}
The analogue of Proposition~5.2 is
\begin{equation}
 \limsup_{\substack{K\to\infty\\40\mid K}}
 \frac{\log m_{9,K,M}}{K^2}
 \leq I_{\mathrm{out},9}+\frac{10\lambda}{M}
                    +\int_3^M\frac{R_9(x)}{x^3}\,dx.
 \label{s5-prime-sum}
\end{equation}
Primes $p\leq\sqrt{5K}$ contribute $O(K^{3/2}\log K)$; this also
absorbs the fixed exceptional factors in $40320$.  For larger primes
in the small range, the logarithmic factor is one, giving
$10\lambda/M$ by the prime number theorem.  For bounded piecewise
continuous $f$, partial summation gives
\[
 K^{-2}\sum_{K/M<p\leq K/3}p f(K/p)\log p
       \longrightarrow\int_3^M f(x)x^{-3}\,dx.
\]
The $O_M(1)$ local errors contribute $o(K^2)$, and the same change of
variables in the outer range gives \eqref{s5-iout9}.

\subsection{The tail integral}
\label{s5-tail}

\paragraph{[Derived] The periodic part is unchanged.}
Put $f=\{x\}$, $g=\{\alpha x\}$, $\tau=\{2Hx\}$,
$\sigma=\{2x\}$, $\eta=\{2\lambda x\}$, and
$A(z)=\ell(x,z)-2x$, $B(z)=\ell(\alpha x,z)-2\alpha x$.
Expanding \eqref{s5-gamma9} gives the analogue of (5.12)--(5.14):
\begin{align}
 R_9(x)&=xF_9(x)+Q(x),\qquad
 F_9(x)=8\lambda+2\lambda f-12\lambda g,\notag\\
 Q(x)&=\frac{\tau(\tau-\sigma)-(\tau-\sigma)_+
                     +\eta(1-\eta)}2
       +9\int_0^{1/2}B^2\,dz-3\int_0^{1/2}AB\,dz,
 \label{s5-tail-decomposition}\\
 -\frac12&\leq Q(x)\leq\frac{13}{8}.\notag
\end{align}
Only the mean linear term has changed.  Define the same periodic primitives
\begin{gather*}
 P(x)=74g(1-g)-\lambda f(1-f),\qquad \overline P=\frac{2923}{240},\\
 G_0(v)=\frac{v(1-v)(2v-1)}6,\qquad
 C_{\mathrm{per}}(x)=\frac{74}{\alpha}G_0(g)-\lambda G_0(f).
\end{gather*}
Then $P'=F_9-3\lambda$, $C_{\mathrm{per}}'=P-\overline P$, and
$|C_{\mathrm{per}}|<16$.  Thus source (5.15) becomes
\begin{align}
 \int_T^\infty\frac{R_9(x)}{x^3}\,dx
 ={}&\frac{3\lambda}{T}-\frac{P(T)}{T^2}
       +\frac{\overline P}{T^2}
       -\frac{2C_{\mathrm{per}}(T)}{T^3}\notag\\
 &+6\int_T^\infty\frac{C_{\mathrm{per}}(x)}{x^4}\,dx
       +\int_T^\infty\frac{Q(x)}{x^3}\,dx.
 \label{s5-tail-identity}
\end{align}
At $T=20$, $P(T)=37/2$ and $C_{\mathrm{per}}(T)=0$.  At $40\mid M$,
both periodic primitives vanish.  Accordingly,
\begin{align}
 \int_{20}^\infty R_9(x)x^{-3}\,dx&\leq\frac{6191}{48000},
 \label{s5-tail-upper}\\
 \int_M^\infty R_9(x)x^{-3}\,dx&\geq\frac{3\lambda}{M}
       +\frac{2923/240-1/4}{M^2}-\frac{32}{M^3}.
 \label{s5-tail-lower}
\end{align}
Appendix~B evaluates the finite integral as
\begin{equation}
 \int_3^{20}R_9(x)x^{-3}\,dx
 =\frac{8222332206149897772829}{7535670527041937280000}.
 \label{s5-finite-inner}
\end{equation}
Combining these expressions yields the analogues of (5.19)--(5.21):
\begin{align}
 A_{9,*}&=\frac{25180312106351244019069}{7535670527041937280000},
 \notag\\
 A_{9,M}&=A_{9,*}+\frac{7\lambda}{M}
             -\frac{2863}{240M^2}+\frac{32}{M^3},
 \label{s5-final-A}\\
 \limsup_{\substack{K\to\infty\\40\mid K}}
 K^{-2}\log m_{9,K,M}&\leq A_{9,M}\qquad(40\mid M).
 \notag
\end{align}
The last inequality is conditional on the local estimates, as before.
The coefficient $7\lambda/M$ survives because the new small-prime cost
$10\lambda/M$ is offset by the tail's $3\lambda/M$.
An especially transparent comparison with the source is
\begin{equation}
 A_{9,M}-A_{5,M}
 =\frac{45541}{57600}
    +4\lambda\left(\frac13-\frac1M\right)+\frac{4\lambda}{M}
 =\frac{116581}{57600}.
 \label{s5-budget-increment}
\end{equation}
All displayed rational integrals are evaluated algebraically.  They
describe the size of this proposed normalization; they do not supply a
missing decay inequality or assert irrationality.
